\begingroup
\renewcommand{\arraystretch}{1.45}
\clearpage\fancyhead[L]{\small ANZEIGE. BESTIMMUNG DER ANZIEHUNG}
\pdfbookmark[0]{ANZEIGE. BESTIMMUNG DER ANZIEHUNG}{work-native366}
\SourceLeaf{366}
\selectlanguage{german}
\begin{center}
\textbf{ANZEIGE.}
\end{center}
\begin{center}
\rule{0.62\linewidth}{0.4pt}\\[6pt]
Göttingische gelehrte Anzeigen.\quad 1818 Februar 9.\\[6pt]
\rule{0.62\linewidth}{0.4pt}
\end{center}

Am $17^{\mathrm{ten}}$ Januar übergab Herr Hofrat \textsc{Gauss} der Königlichen Gesellschaft eine Vorlesung:

\begin{center}
\textit{Determinatio attractionis, quam in punctum quodlibet positionis datae exerceret
planeta, cuius massa per totam eius orbitam, ratione temporis, quo singulae partes
describuntur, uniformiter esset dispertita.}
\end{center}

\Needspace{27\baselineskip}
Nach einem Lehrsatz, den vielleicht bisher noch niemand ausdrücklich ausgesprochen hat, der sich aber aus den Grundlagen der physikalischen Astronomie leicht beweisen lässt, sind die säkularen Veränderungen einer Planetenbahn durch die Störung eines anderen Planeten dieselben, ob der störende Planet seine elliptische Bahn nach Keplers Gesetzen tatsächlich durchläuft oder ob man annimmt, seine Masse sei auf den Umfang der Ellipse so verteilt, dass auf Stücke der Ellipse, die sonst in gleich großen Zeiten durchlaufen werden. gleich große Anteile der gesamten Masse entfallen: vorausgesetzt, dass die Umlaufzeiten des gestörten und des störenden Planeten nicht in einem rationalen Verhältnis zueinander stehen. Die Aufgabe, die den Gegenstand dieser Abhandlung bildet, nämlich die nicht angenäherte, sondern genaue, nicht von mäßiger Exzentrizität der Ellipse abhängige. sondern allgemeine Bestimmung der Anziehung, die ein elliptischer Ring von unendlich kleiner und nach obigem
\SourceLeaf{367}
Gesetz unveränderlicher Dicke auf jeden Punkt im Raum ausübt, ist daher für die physikalische Astronomie von hohem Interesse. Sie ist aber auch in rein mathematischer Hinsicht nicht weniger bemerkenswert, wegen der verschiedenen Kunstgriffe, die ihre vollständige Lösung erfordert.

\Needspace{40\baselineskip}
Von der Lösung selbst lässt sich hier nicht gut ein Auszug geben. Der Verfasser hat eine rein analytische Behandlung gewählt; Kenner, die ihr aufmerksam folgen, werden leicht die geometrischen Entsprechungen der einzelnen in der Untersuchung vorkommenden Größen erkennen und sehen, wie sie sich in eine geometrische Form umgießen lässt. Hier möge es genügen, nur das Endergebnis anzuführen. Drei unbekannte Größen $G,\ G',\ G''$ werden durch die Wurzeln einer kubischen Gleichung bestimmt, aus deren Beschaffenheit sich beweisen lässt, dass sie stets drei reelle Wurzeln hat. Die nach einer beliebigen Richtung zerlegte Anziehung des elliptischen Ringes wird dann durch einen Ausdruck der Form $pP+qQ$ dargestellt, wobei $p$ und $q$ algebraisch von $G,\ G',\ G''$ abhängen, während $P$ und $Q$ die bestimmten Werte der Integrale
\[
\int\frac{\cos T^{2}\,.\,\mathrm{d}T}
{2\pi(mm\cos T^{2}+nn\sin T^{2})^{\frac{3}{2}}}
\]
\[
\int\frac{\sin T^{2}\,.\,\mathrm{d}T}
{2\pi(mm\cos T^{2}+nn\sin T^{2})^{\frac{3}{2}}}
\]
bedeuten, wenn die Integrationen von $T=0$ bis $T=360^{0}=2\pi$ erstreckt werden, und wobei
\[
m=\surd(G+G'),\quad n=\surd(G+G'')
\]
\Needspace{45\baselineskip}
\par Da diese Integrale (ausgenommen $m=n$) transzendenter Natur sind. und bekanntermaßen mit anderen viel behandelten transzendenten Größen der Störungsrechnung zusammenhängen, konnte die Lösung, nachdem sie bis zu diesem Punkt geführt war, als vollendet angesehen werden. Der Verfasser hat jedoch diese erste sich ihm bietende Gelegenheit benutzt, um die ersten Grundzüge eines \textit{neuen Algorithmus} anzugeben, den er schon seit vielen Jahren zur Bestimmung dieser transzendenten Größen benutzt hat und über den er künftig eine ausführliche, zu vielen bemerkenswerten Ergebnissen führende Untersuchung veröffentlichen wird. Hier können nur die Hauptsätze unter Weglassung der Beweise angeführt werden. Leitet man aus zwei gegebe\SourceLeaf{368}%
nen positiven Größen $m$ und $n$ weitere Größen $m',\ m'',\ m'''$ usw., $n',\ n'',\ n'''$ usw. nach folgenden Regeln ab:
\[
\begin{array}{ll}
m'=\frac{1}{2}(m+n),&n'=\surd mn\\[0.4em]
m''=\frac{1}{2}(m'+n'),&n''=\surd m'n'\\[0.4em]
m'''=\frac{1}{2}(m''+n''),&n'''=\surd m''n''\ \mbox{usw.}
\end{array}
\]
\Needspace{27\baselineskip}
das heißt, sind $m',\ n'$ das arithmetische beziehungsweise geometrische Mittel zwischen $m$ und $n$ und ebenso $m'',\ n''$ das arithmetische und geometrische Mittel zwischen $m'$ und $n'$ und so fort, so nähern sich die Glieder sowohl der Folge $m,\ m',\ m'',\ m'''$ usw. als auch der Folge $n,\ n',\ n'',\ n'''$ usw. äußerst schnell einem gemeinsamen Grenzwert $=\mu$, den der Verfasser das arithmetisch-geometrische Mittel von $m$ und $n$ nennt. Offenbar ist $\mu$ zugleich das arithmetisch-geometrische Mittel von $m'$ und $n'$ oder überhaupt von je zwei zusammengehörigen Gliedern der beiden Folgen. Der Verfasser beweist nun, dass $\frac{1}{\mu}$ der Wert des Integrals
\[
\int\frac{\mathrm{d}T}{2\pi\surd(mm\cos T^{2}+nn\sin T^{2})}
\]
ist, wenn die Integration von $T=0$ bis $T=360^{0}$ erstreckt wird. Man wird leicht sehen, dass sich dies auch folgendermaßen hätte aussprechen lassen: Ergibt die Entwicklung der Funktion
\[
\frac{1}{\surd(\alpha+\beta\cos\psi)}
\]
die Reihe
\[
A+B\cos\psi+C\cos2\psi+D\cos3\psi+\mathrm{usw.}
\]
so ist $\frac{1}{A}$ stets das arithmetisch-geometrische Mittel der beiden Größen $\surd(\alpha+\nobreak{}\beta)$ und $\surd(\alpha-\nobreak{}\beta)$.

\Needspace{38\baselineskip}
Ein zweiter, ebenso wichtiger Satz besagt, dass, wenn man die Summe der unendlichen, stets sehr schnell konvergierenden Reihe
\[
2(m'm'-n'n')+4(m''m''-n''n'')+8(m'''m'''-n'''n''')+\mbox{usw.}
\]
bei der die Zahlenkoeffizienten eine geometrische Folge bilden, $=(mm-nn)\nu$ setzt, der Wert des Integrals
\[
\int\frac{\cos2T\,.\,\mathrm{d}T}
{2\pi\surd(mm\cos T^{2}+nn\sin T^{2})}
\]
\SourceLeaf{369}
von $T=0$ bis $T=360^{0}$ erstreckt $=-\frac{\nu}{\mu}$ wird. Offenbar ist hierdurch auch der zweite Koeffizient der obigen Reihe bekannt, nämlich $B=-\frac{\nu}{2\mu}$, wenn man $m=\surd(\alpha+\beta)$, $n=\surd(\alpha-\beta)$ gesetzt hat. Alle folgenden Koeffizienten $C$, $D$ usw. werden jedoch bekanntlich durch die beiden ersten $A$ und $B$ algebraisch und einfach bestimmt. Für die numerische Berechnung der Größen $m'm'-n'n'$, $m''m''-n''n''$ usw. wird in der Abhandlung selbst noch ein besonderes, sehr bequemes Verfahren angegeben.
\par
Die Anwendung auf die transzendenten Größen der vorliegenden Untersuchung ergibt schließlich noch die einfachen Ausdrücke
\[
P=\frac{1+\nu}{2mm\mu},\qquad Q=\frac{1-\nu}{2nn\mu}
\]
\Needspace{32\baselineskip}
\par Aufmerksamen Lesern wird nicht entgehen, wie viele interessante Aufgaben, die mit den hier betrachteten transzendenten Größen zusammenhängen, durch den erläuterten Algorithmus mit größter Leichtigkeit gelöst werden. Als Beispiel führen wir hier die Rektifikation der Ellipse an. Setzt man ihre halbe große Achse $=m$, ihre halbe kleine Achse $=n$, so wird der Umfang
\[
\begin{aligned}&=\frac{2\pi}{\mu}\bigl\{m'm'-2(m''m''-n''n'')\bigr.\\[5pt]&\qquad\bigl.-4(m'''m'''-n'''n''')-8(m''''m''''-n''''n'''')-\mathrm{usw.}\bigr\}\end{aligned}
\]
Ein anderes Beispiel ist die Dauer der Pendelschwingungen bei endlichen Bögen; sie verhält sich zur Dauer der unendlich kleinen Schwingungen wie die Einheit zum arithmetisch-geometrischen Mittel zwischen $1$ und dem Kosinus eines Viertels des gesamten Schwingungsbogens.
\Needspace{14\baselineskip}
\par
Schließlich muss noch bemerkt werden, dass der Verfasser glaubte, diese Ergebnisse in ihrer ursprünglichen Form darstellen zu müssen, so wie er sie schon vor vielen Jahren unabhängig von ähnlichen Untersuchungen \textsc{Lagranges} und \textsc{Legendres} gefunden hatte, obwohl sie zum Teil leicht aus den Entdeckungen dieser Geometer hätten abgeleitet werden können: teils weil jene Form ihm wesentliche Vorzüge zu haben schien, teils weil sie gerade so den Anfang einer viel ausgedehnteren Theorie bilden, in der seine Arbeit eine ganz andere Richtung genommen hat als die der genannten Geometer.
\begin{center}\rule{0.25\linewidth}{0.4pt}\end{center}
\clearpage
\endgroup
